% =====================================================================
%  Curriculum Vitae -- Newton Huy Nguyen
%  Build:  pdflatex CV_Newton.tex   (run once; no bibliography needed)
%  Output: CV_Newton.pdf  -> copy to ../pdfs/CV_Newton.pdf
%
%  SOURCES MERGED INTO THIS VERSION (Nov 2026):
%    - CV_Newton.pdf as uploaded to GitHub, Apr 2024  (academic sections)
%    - Newton_Nguyen_Spectral_ML_Resume-2.docx, 2026  (positions, systems)
%    - OpenAlex publication record                    (verified citations)
%
%  STILL TO FILL IN -- only you know these:
%    - Honors or awards received in 2025 and 2026
%    - Invited talks given after the Stanford d.School panel (2023)
%    - Conference presentations after AGU 2022 that you presented yourself
%      (the 2023/2025 Optica and 2026 EGU abstracts are listed under
%       publications; move them here if you presented them)
%    - Teaching after Summer 2019, including the Caltech co-taught courses
%    - End dates for service roles still marked "Present"
% =====================================================================

\documentclass[11pt,letterpaper]{article}

\usepackage[utf8]{inputenc}
\usepackage[T1]{fontenc}
\usepackage{newtxtext}
\usepackage[margin=0.8in,bottom=0.75in]{geometry}
\usepackage{enumitem}
\usepackage{titlesec}
\usepackage{xcolor}
\usepackage[hidelinks]{hyperref}
\usepackage{fancyhdr}
\usepackage{needspace}

\definecolor{rule}{HTML}{30363D}
\definecolor{link}{HTML}{1D6670}
\hypersetup{colorlinks=true, urlcolor=link, linkcolor=link,
            pdftitle={Newton Huy Nguyen -- Curriculum Vitae},
            pdfauthor={Newton Huy Nguyen}}

% --- section headings -------------------------------------------------
\titleformat{\section}
  {\normalfont\bfseries\scshape\large\color{rule}}
  {}{0pt}{}[{\color{rule}\titlerule[0.9pt]}]
\titlespacing*{\section}{0pt}{11pt}{5pt}

\titleformat{\subsection}{\normalfont\bfseries\itshape\normalsize}{}{0pt}{}
\titlespacing*{\subsection}{0pt}{9pt}{3pt}

% --- helpers ----------------------------------------------------------
\setlist[itemize]{leftmargin=1.3em, topsep=2pt, itemsep=1pt, parsep=0pt}
\newcommand{\entry}[3]{\noindent\textbf{#1}~\textbar~#2\hfill{#3}\par}
\newcommand{\simpleentry}[2]{\noindent #1\hfill{#2}\par\vspace{1pt}}
\newcommand{\me}{\textbf{N.H. Nguyen}}
\setlength{\parindent}{0pt}

\newenvironment{pubs}
  {\begin{list}{}{\leftmargin=1.6em \itemindent=-1.6em
   \topsep=3pt \itemsep=4pt \parsep=0pt}}
  {\end{list}}

% keep a section heading with at least a few lines of its content
\let\oldsection\section
\renewcommand{\section}[1]{\Needspace*{5\baselineskip}\oldsection{#1}}
\let\oldsubsection\subsection
\renewcommand{\subsection}[1]{\Needspace*{4\baselineskip}\oldsubsection{#1}}

\pagestyle{fancy}
\fancyhf{}
\renewcommand{\headrulewidth}{0pt}
\fancyfoot[C]{\small\color{gray}Newton Huy Nguyen \textbullet\ \thepage}

\begin{document}

% =====================================================================
\begin{center}
  {\LARGE\bfseries Newton Huy Nguyen, Ph.D.}\\[4pt]
  Doerr School of Sustainability \& Department of Physics, Stanford University\\[3pt]
  San Francisco, CA \quad\textbullet\quad (408) 613-4379 \quad\textbullet\quad
  \href{mailto:nnewton@stanford.edu}{nnewton@stanford.edu}\\[2pt]
  \href{https://newton-climate.github.io}{newton-climate.github.io} \;\textbullet\;
  \href{https://scholar.google.com/citations?user=AziOzdwAAAAJ&hl=en}{Google Scholar} \;\textbullet\;
  \href{https://orcid.org/0000-0002-9118-8672}{ORCID 0000-0002-9118-8672} \;\textbullet\;
  \href{https://github.com/Newton-Climate}{GitHub}
\end{center}
\vspace{4pt}

% =====================================================================
\section{Education}

\entry{Ph.D., California Institute of Technology}{Environmental Engineering Science}{2019--2023}
\textit{Thesis:} From Source to Sink: Measuring and Modeling Methane Emissions and Loss.\\
\textit{Advisor:} Christian Frankenberg
\vspace{5pt}

\entry{M.S., California Institute of Technology}{Environmental Engineering Science}{2017--2019}
\textit{Thesis:} Effects of Coupled Chemistry on Methane Emissions Estimates.\\
\textit{Advisors:} Tapio Schneider \& Christian Frankenberg
\vspace{5pt}

\entry{B.A., University of California, Berkeley}{Geophysics, Highest Honors}{2012--2016}
\textit{Thesis:} Neural Networks to Model Fluid Flows.\\
\textit{Advisor:} Bruce Buffett

% =====================================================================
\section{Positions}

\entry{Stanford Science Fellow, Postdoctoral Scholar}{Stanford University}{Sep 2023--Sep 2024; Feb 2026--Present}
\textit{Supervisors:} Alison Hoyt \& Leo Hollberg
\begin{itemize}
  \item Co-leading a group developing ecosystem-scale greenhouse gas measurements with long-range lasers.
  \item Quantifying the sensitivity of soil carbon sequestration to environmental change.
  \item Built a retrieval algorithm for Earth's surface properties from far-infrared satellite spectroscopy, for processing measurements from the NASA PREFIRE mission.
  \item Designed observing-system and data-assimilation simulations to optimize sensor placement and quantify information gain under physical and measurement uncertainty.
  \item Developed hybrid machine-learning and physical models, with Bayesian information-content diagnostics, to identify which observations best constrain model parameters.
  \item Mentoring a Ph.D. researcher in remote sensing, model development, data assimilation, and production-quality scientific code.
\end{itemize}
\vspace{4pt}

\entry{Atmospheric Data Scientist; Data \& Algorithms Consultant}{Puloli, Inc.}{Jan 2023--Feb 2026}
Consultant Jan 2023--Sep 2024; full-time data scientist Sep 2024--Feb 2026.
\begin{itemize}
  \item Algorithm and data-pipeline development for continuous methane monitoring in oil and gas.
  \item Built Python analysis and scraping pipelines integrating surface and high-altitude balloon measurements into asynchronous deployment workflows.
  \item Designed a real-time model-quality metric quantifying forecast uncertainty for operational decisions, increasing accuracy and automating measurement decisions (patent being filed).
\end{itemize}
\vspace{4pt}

\entry{Algorithms \& Scientific Computing Consultant}{Multiple clients}{2023--Present}
\begin{itemize}
  \item Designed and deployed an AWS observation-planning system for a hyperspectral satellite constellation, reducing runtime from 46 hours to 4 while optimizing sensing opportunities.
  \item Led customer workshops translating sensor and algorithm needs into technical requirements, validation protocols, and end-to-end tests.
\end{itemize}
\vspace{4pt}

\entry{Ph.D. Researcher}{Caltech / NASA Jet Propulsion Laboratory}{Sep 2017--Jul 2023}
\textit{Advisor:} Christian Frankenberg
\begin{itemize}
  \item Built Bayesian inversion and 1-D spectral-retrieval algorithms for greenhouse gases across satellite, Fourier-transform infrared, dual-comb, and laser-heterodyne measurements; accelerated inference on GPU and parallel systems.
  \item Awarded an NSF Graduate Research Fellowship to modernize methane monitoring capabilities.
  \item Mentored two students and one software engineer in satellite remote sensing techniques.
\end{itemize}
\vspace{4pt}

\entry{Research Assistant}{Lawrence Berkeley National Laboratory}{Jun 2016--Jul 2017}
\textit{Supervisors:} Daniel Feldman \& William D. Collins
\begin{itemize}
  \item Team member on NASA's CLARREO Science Team; investigated cloud--climate feedback using NASA satellite products.
  \item Parallelized a PCA-based radiative transfer model on a 128-node cluster (MPI and Fortran).
\end{itemize}
\vspace{4pt}

\entry{Summer Research Fellow}{Computational \& Applied Sciences Lab, UC Santa Barbara}{Summer 2015}
\textit{Advisor:} Fr\'ed\'eric Gibou
\vspace{3pt}

\entry{Summer Research Intern}{Hydro-seismology Lab, UC Berkeley}{Summer 2014}
\textit{Advisor:} Chi Wang

% =====================================================================
\section{Honors \& Awards}

\simpleentry{Stanford University Science Fellowship}{2023--2026}
\simpleentry{Caltech Engineering Division New Horizons Prize for Excellence in Mentorship \& Service}{2023}
\simpleentry{Research University Alliance}{2022}
\simpleentry{National Science Foundation Graduate Research Fellowship}{2018}
\simpleentry{Boston Marathon Qualifier}{2019, 2020}
\simpleentry{3rd Place, US Blind Athletes National Championships in the Marathon}{2019}
\simpleentry{1st Place, Collegiate Triathlon National Championships, para-athlete division}{2016}
\simpleentry{``UC Leads'' Fellowship \& Best Presentation Award}{2014--2016}
\simpleentry{Albert Newman Fellowship}{2014}
\simpleentry{Berkeley Academic Merit Scholarship}{2013--2016}

% =====================================================================
\section{Publications}

\subsection{Peer-reviewed journal articles}
\begin{pubs}
\item[6.] K.A. Lamantia, L.G. Thompson, M.E. Davis, B. Riddell-Young, I. Strawson, É. Beaudon,
  E. Mosley-Thompson, \me, E.J. Brook (2026). A global atmospheric methane record from a tropical
  ice core. \textit{Nature}. \href{https://doi.org/10.1038/s41586-026-10938-1}{10.1038/s41586-026-10938-1}

\item[5.] H. Dion-Kirschner, \me, C. Frankenberg, W.W. Fischer (2024). Evaluating the contribution of
  methanotrophy kinetics to uncertainty in the soil methane sink. \textit{Environmental Research Letters}.
  \href{https://doi.org/10.1088/1748-9326/ad4c7a}{10.1088/1748-9326/ad4c7a}

\item[4.] R.K. Cole, C.D. Fredrick, \me, S.A. Diddams (2023). Precision Doppler shift measurements with a
  frequency-comb-calibrated laser heterodyne radiometer. \textit{Optics Letters}.
  \href{https://doi.org/10.1364/OL.500652}{10.1364/OL.500652}

\item[3.] \me, A.J. Turner, Y. Yin, M.J. Prather, C. Frankenberg (2020). Effects of chemical feedbacks on
  decadal methane emissions estimates. \textit{Geophysical Research Letters}.
  \href{https://doi.org/10.1029/2019GL085706}{10.1029/2019GL085706}

\item[2.] W.D. Collins, D.R. Feldman, C. Kuo, \me\ (2018). Large regional shortwave forcing by
  anthropogenic methane informed by Jovian observations. \textit{Science Advances}.
  \href{https://doi.org/10.1126/sciadv.aas9593}{10.1126/sciadv.aas9593}

\item[1.] D.R. Feldman, W.D. Collins, Y.L. Shea, \me, X. Liu, B.A. Wielicki (2016). Observing climate change
  with both shortwave and longwave hyperspectral satellite instrumentation. \textit{Light, Energy and the
  Environment}. \href{https://doi.org/10.1364/HISE.2016.HW2F.1}{10.1364/HISE.2016.HW2F.1}
\end{pubs}

\subsection{Conference proceedings \& abstracts}
\begin{pubs}
\item[3.] Y. Yin, \me, Q. Wang, A.J. Turner, F. Chevallier, C. Frankenberg (2026). Natural variability in
  stratosphere--troposphere exchange: an underappreciated driver of global methane growth rates.
  \textit{EGU General Assembly}.
  \href{https://doi.org/10.5194/egusphere-egu26-21945}{10.5194/egusphere-egu26-21945}

\item[2.] C. Huff, A.M. Hoyt, J. Lamb, \me, A. Majumdar, L.W. Hollberg (2025). Enhancing the monitoring of
  dispersed methane sources with open-path optical sensors. \textit{Optica Sensing Congress}.
  \href{https://doi.org/10.1364/SENSORS.2025.SM2D.2}{10.1364/SENSORS.2025.SM2D.2}

\item[1.] R.K. Cole, C.D. Fredrick, \me, C. Frankenberg, S.A. Diddams (2023). Precision atmospheric wind
  measurements with a frequency-comb-calibrated laser heterodyne radiometer.
  \textit{Optica Sensing Congress}.
  \href{https://doi.org/10.1364/ES.2023.ETu5E.5}{10.1364/ES.2023.ETu5E.5}
\end{pubs}

\subsection{Submitted or in preparation}
\begin{pubs}
\item[3.] \me, K. Cossel, E. Waxman, N. Newbury, I. Coddington, C. Frankenberg (submitted). Towards
  laboratory-level accuracy in the field: environmental impacts on greenhouse gas observations and
  spectroscopy measured by laser frequency combs. \textit{Atmospheric Measurement Techniques}.

\item[2.] \me, K. Cossel, E. Waxman, N. Newbury, I. Coddington, C. Frankenberg (in prep).
  Next-generation ground-based measurements of vertical greenhouse gas gradients with laser frequency
  combs. \textit{Atmospheric Measurement Techniques}.

\item[1.] \me, Y. Yin, A.J. Turner, C. Frankenberg (in prep). Global methane emissions trends and
  atmospheric oxidation are modulated by stratospheric transport and ENSO.
\end{pubs}

\subsection{Other publications}
\begin{pubs}
\item[3.] D.R. Dounas-Frazer et al.\ (2022). Increase investment in accessible physics labs: a call to action
  for the physics education community. \textit{arXiv preprint}, endorsed by the American Association of
  Physics Teachers. \href{https://doi.org/10.48550/arXiv.2202.00816}{10.48550/arXiv.2202.00816}

\item[2.] Caltech (2021). \href{https://inclusive.caltech.edu/documents/19303/Report_of_the_Committee_on_Student_Admissions_and_Recruitment.pdf}{Report
  of the Committee on Student Admissions and Recruitment}.

\item[1.] Glassdoor (2017). \href{https://www.glassdoor.com/blog/newton-nguyen/}{A triathlete shares his
  5 secrets to success}.
\end{pubs}

% =====================================================================
\section{Research Software}

\entry{EnforceFlux}{Python}{2024--Present}
Observation-simulator framework for simulating heterogeneous greenhouse gas sources and evaluating
monitoring network layouts across atmospheric transport models (FLEXPART).
\href{https://github.com/Newton-Climate/EnforceFlux}{github.com/Newton-Climate/EnforceFlux}
\vspace{4pt}

\entry{ConstrainedCarbon}{Python}{2024--Present}
Bayesian observing-system simulator using optimal estimation, degrees of freedom for signal, and Shapley
attribution to rank the value of candidate measurements.
\href{https://github.com/Newton-Climate/ConstrainedCarbon}{github.com/Newton-Climate/ConstrainedCarbon}
\vspace{4pt}

\entry{SpectralFit.jl}{Julia}{Jun 2020--Present}
Open-source Bayesian retrieval engine for greenhouse gas concentrations and vertical profiles from TCCON,
OCO-2, and dual-comb spectra, with GPU and parallel acceleration. Resulted in two invited talks and one
peer-reviewed publication.
\href{https://github.com/Newton-Climate/SpectralFit.jl}{github.com/Newton-Climate/SpectralFit.jl}
\vspace{4pt}

\entry{OHMethane}{Julia \& MATLAB}{Jan 2018--Present}
Four-box model simulating atmospheric methane chemistry and inferring global emissions under chemical
constraints via Bayesian optimization. Resulted in two conference presentations and two peer-reviewed
publications. \href{https://github.com/Newton-Climate/OHMethane}{github.com/Newton-Climate/OHMethane}

% =====================================================================
\section{Invited Presentations}

\subsection{Technical talks}
\simpleentry{University of Washington, Atmospheric Science Dept.\ Seminar \textbar\ ``Towards continuous, 3-D observation of greenhouse gas concentrations directly in the field''}{2022}
\simpleentry{University of Texas at Austin, Dept.\ Seminar \textbar\ ``From source to sink: constraining past \& present methane emissions''}{2022}
\simpleentry{Harvard, Atmospheric Chemistry Group \textbar\ ``Biological, dynamic, and chemical drivers of methane destruction''}{2021}
\simpleentry{Caltech, Environmental Science Dept.\ Seminar \textbar\ ``Destructive interference: frequency combs for greenhouse gas remote sensing''}{2021}
\simpleentry{Caltech, Astronomy Division \textbar\ ``Monitoring climate change from space''}{2019}

\subsection{Invited outreach talks}
\simpleentry{Stanford d.school \textbar\ ``Panel on accessibility for engineering education''}{2023}
\simpleentry{Coca-Cola Headquarters \textbar\ ``Activism and perseverance in the Asian community''}{2022}
\simpleentry{Reclaiming STEM \textbar\ ``Disability in science''}{2021}
\simpleentry{Communication Science Conference (ComSciCon) \textbar\ ``DEI in STEM''}{2021}
\simpleentry{National Association of Blind Students \textbar\ ``Communicating as a blind presenter''}{2021}
\simpleentry{National Federation of the Blind \textbar\ ``Designing your environment: how to succeed as a blind scientist''}{2021}
\simpleentry{National Federation of the Blind \textbar\ ``Developing tools for blind scientists''}{2019}
\simpleentry{Aira Corporation \textbar\ ``Being a blind scientist''}{2018}

% =====================================================================
\section{Conference Presentations}

\begin{pubs}
\item[] \me, K. Cossel, E. Waxman, N. Newbury, I. Coddington, C. Frankenberg (2022). Bridging the gap:
  frequency combs for a next-generation GHG monitoring network. Poster, American Geophysical Union.

\item[] \me, K. Cossel, E. Waxman, N. Newbury, I. Coddington, C. Frankenberg (2021). Destructive
  interference: future long-term greenhouse gas monitoring with dual-comb spectroscopy needs more
  accurate spectroscopic parameters. Poster, American Geophysical Union.

\item[] \me, C. Frankenberg, A.J. Turner, Y. Yin, M.J. Prather (2019). Quantifying the effect of neglecting
  variable methane lifetime on methane emissions estimates. Poster, American Geophysical Union.

\item[] \me, C. Frankenberg, Y. Yin, A.J. Turner (2018). Effects of methane and hydroxyl radical chemistry
  on decadal methane emissions estimates. Poster, American Geophysical Union.

\item[] D. Feldman, W.D. Collins, B.A. Wielicki, Y. Shea, M.G. Mlynczak, C. Kuo, \me\ (2017). How
  continuous observations of shortwave reflectance spectra can narrow the range of shortwave climate
  feedbacks. Poster, American Geophysical Union.
\end{pubs}

% =====================================================================
\section{Teaching}

\simpleentry{Caltech \textbar\ Co-taught courses in infrared sensing, numerical modelling of physical systems, and information retrieval theory}{2019--2023}
\simpleentry{ESE 103 Biogeochemistry \textbar\ Teaching assistant, Christian Frankenberg (rating 5/5)}{Spring 2019}
\simpleentry{ESE 156 Remote Sensing of the Atmosphere \& Biosphere \textbar\ Teaching assistant, Christian Frankenberg (rating 5/5)}{Fall 2019}
\simpleentry{Scientific Writing \& Communication, UCLA \textbar\ Teaching assistant}{Summer 2019}

% =====================================================================
\section{Leadership \& Service}

\entry{Co-founder}{Systemic Access Mentorship Program}{Aug 2020--Present}
\begin{itemize}
  \item National mentorship program for blind students in STEM, with 40 participants across three countries.
  \item Convene virtual meetings twice a month and meet with mentees regularly.
\end{itemize}
\vspace{3pt}

\simpleentry{\textbf{Co-founder, President \& Treasurer}~\textbar~Caltech Disability Coalition}{2020--2022}
\vspace{3pt}

\entry{Student representative}{Caltech Graduate Admissions Policy Committee}{Sep 2020--Mar 2021}
\begin{itemize}
  \item Selected by the President of Caltech for the faculty committee tasked with increasing student body diversity and rewriting graduate admissions policies.
  \item Proposed three policies for reducing admissions bias, which the university implemented.
\end{itemize}
\vspace{3pt}

\entry{Member}{AAPT Committee for Accessible Labs}{2019--2021}
\begin{itemize}
  \item Improving disability access in K--12 and university physics labs; produced a white paper on best practices and a conference presentation.
\end{itemize}
\vspace{3pt}

\entry{Representative, Engineering \& Applied Sciences}{Caltech Graduate Student Council}{2018--2020}
\begin{itemize}
  \item Member of the Advocacy \& Diversity Sub-committee.
\end{itemize}
\vspace{3pt}

\entry{Founder \& President}{Caltech Triathlon Club}{Sep 2019--Present}
\begin{itemize}
  \item Coached track and cycling practices for more than 20 athletes.
  \item Organized a three-race series involving 100 participants from 11 schools.
\end{itemize}
\vspace{3pt}

\entry{Committee member}{Lawrence Berkeley National Laboratory}{2016--2017}
\begin{itemize}
  \item Lab-wide Employee Accessibility Committee; Climate \& Ecosystem Science Division representative for diversity and inclusion; DEI representative for the DOE External Review Committee.
\end{itemize}
\vspace{3pt}

\simpleentry{\textbf{President}~\textbar~UC Berkeley Atmospheric Science Association}{2015--2016}

% =====================================================================
\section{Selected Press}

\simpleentry{Stanford University \textbar\ 2023 cohort of Stanford Science Fellows named}{2023}
\simpleentry{\#SoCaltech \textbar\ Newton Nguyen}{Mar 2022}
\simpleentry{Runner's World \textbar\ How running has helped Newton Nguyen navigate life as a blind man}{Oct 2021}
\simpleentry{Ally commercial \textbar\ We're all better off with an Ally}{Jul 2021}
\simpleentry{Mini-documentary \textbar\ Marathoning through life with vision loss}{Oct 2021}
\simpleentry{Careers for the Blind Podcast \textbar\ Episode 35: Climate Scientist Newton Nguyen}{Oct 2021}
\simpleentry{Caltech Magazine \textbar\ Creating a more inclusive Caltech}{Fall 2020}
\simpleentry{ScienceDaily \textbar\ By Jove! Methane's effects on sunlight vary by region}{Sep 2018}
\simpleentry{Berkeley Lab News \textbar\ Newton Nguyen's vision loss doesn't slow him down}{Oct 2016}
\simpleentry{KTVU News \textbar\ 98 percent blind Cal Berkeley student competes as triathlete}{Mar 2015}

% =====================================================================
\section{Technical Skills}

\textbf{Languages \& tools:} Python, Julia (incl.\ CUDA), R, MATLAB, Fortran, MPI, Git, Bash, Docker, AWS,
GitHub Actions, AsyncIO.\\[2pt]
\textbf{Scientific stack:} NumPy/SciPy, PyTorch, JAX, scikit-learn, Xarray, NetCDF, Zarr, HDF5.\\[2pt]
\textbf{Expertise:} Bayesian inference and optimal estimation, uncertainty quantification, data assimilation,
observation planning and information content, radiative transfer, satellite and ground-based spectroscopy,
greenhouse gas emissions, high-performance and GPU computing, machine learning.\\[2pt]
\textbf{Memberships:} American Geophysical Union (2014--present); American Meteorological Society
(2017--present).

\end{document}
